% ============================================================================
% Fichier  : partie6/P06_C31_methodologie_formelle.tex
% Titre    : Méthodologie d'Analyse Formelle de Protocoles
% Auteur   : MINKA MI NGUIDJOÏ Thierry Emmanuel
% Version  : 2.0 -- Juin 2026
% ============================================================================

\chapter{Méthodologie d'Analyse Formelle de Protocoles}
\label{chap:form-meth}

\epigraph{%
    « Sans modélisation formelle, la sécurité n'est qu'une illusion de
    confiance. »%
}{Andrew Yao}

\niveauChapitre{Recherche}{%
    Logique formelle de base recommandée. Chapitre~\ref{chap:cryptanal}
    (cryptanalyse) indispensable~: ce chapitre formalise et outille
    la démarche cryptanalytique qui y a été présentée.%
}

\begin{boxobjectifs}
\begin{itemize}
    \item Maîtriser le modèle de menace Dolev-Yao~\cite{dolevyao1983}~: ses hypothèses,
          sa portée, et pourquoi il reste le standard de l'analyse
          de protocoles malgré ses limites assumées.
    \item Formaliser des propriétés de sécurité (non-répudiation,
          confidentialité, authenticité, intégrité) sous forme de
          lemmes vérifiables.
    \item Comprendre les principes de fonctionnement et les usages
          respectifs de Tamarin Prover, ProVerif et CryptoVerif.
    \item Appliquer une méthodologie d'audit en cinq étapes,
          combinant analyse manuelle et vérification automatisée.
    \item Interpréter correctement un résultat ``non prouvé'' rendu
          par un outil de vérification formelle~: distinguer
          limitation outil et faille réelle.
    \item Relier la vérification formelle de protocoles à la
          méthodologie de rédaction de rapport forensique
          (Ch.~\ref{chap:rapport})~: même exigence de rigueur,
          même distinction entre ce qui est prouvé et ce qui ne
          l'est pas.
\end{itemize}
\end{boxobjectifs}

\bigskip

% ============================================================================
\section*{De la cryptanalyse informelle à la preuve formelle}
% ============================================================================

Le Chapitre~\ref{chap:cryptanal} a présenté la cryptanalyse comme
discipline de test adverse, indispensable au même titre que la
falsification active du PIU (E9bis, Ch.~6~(UE~SEC4031)). Mais
``tester un protocole'' et ``prouver qu'il résiste à toute attaque
dans un modèle donné'' sont deux démarches de nature différente~:
la première est empirique et jamais exhaustive~; la seconde, quand
elle aboutit, offre une garantie mathématique au sein des hypothèses
posées. Ce chapitre présente la méthodologie et les outils qui
permettent de passer de l'une à l'autre.

\separateur

% ============================================================================
\section{Modélisation des Menaces}
\label{sec:modelisation-menaces}
% ============================================================================

Toute analyse formelle commence par la définition précise de la
puissance et des objectifs de l'adversaire considéré. Sans cette
définition explicite, affirmer qu'un protocole est ``sûr'' n'a
aucun sens vérifiable, sûr contre quel adversaire, capable de
quoi~?

\subsection{Le modèle Dolev-Yao}

\begin{boxdef}
\textbf{Modèle Dolev-Yao}~(1983)~: modèle standard pour l'analyse
formelle des protocoles cryptographiques. Il pose que l'attaquant~:
\begin{itemize}
    \item \textbf{contrôle entièrement le réseau}~: il peut écouter,
          bloquer, injecter et modifier n'importe quel message en
          transit~;
    \item \textbf{est un participant légitime}~: il possède les
          clés publiques attendues et peut engager des sessions
          comme n'importe quel acteur honnête~;
    \item \textbf{ne peut pas casser les primitives cryptographiques
          par la force}~: le chiffrement, les signatures et les
          fonctions de hachage sont traités comme des
          \textbf{boîtes noires parfaites} (modèle de l'oracle), 
          un message chiffré est totalement opaque sans la clé,
          une signature ne peut être forgée sans la clé privée.
\end{itemize}
\end{boxdef}

\begin{boxalert}
\textbf{Pourquoi cette séparation des préoccupations est une force,
pas une faiblesse}

Le modèle Dolev-Yao \textbf{ignore délibérément} la cryptanalyse des
primitives (objet du Ch.~\ref{chap:cryptanal}, Section~``contre les
primitives'' du Tableau~30.1) pour se concentrer exclusivement sur
les \textbf{failles logiques du protocole}~: l'agencement des
messages, l'ordre des étapes, la gestion des identités. Cette
séparation permet d'isoler deux questions distinctes~: ``la
mathématique sous-jacente est-elle solide~?'' (cryptanalyse des
primitives, Ch.~\ref{chap:cryptanal}) et ``même avec une mathématique
parfaite, le protocole est-il logiquement bien conçu~?'' (analyse
formelle, ce chapitre). Un protocole peut échouer sur l'une de ces
deux dimensions indépendamment de l'autre, la taxonomie du
Tableau~30.1 distingue précisément les failles ``contre le
protocole'' (logiques, du ressort de Dolev-Yao) des failles ``contre
les primitives'' (mathématiques, hors du périmètre Dolev-Yao).
\end{boxalert}

\subsection{Formalisation des propriétés de sécurité}

Les propriétés de sécurité attendues d'un protocole doivent être
exprimées de manière formelle et mécaniquement vérifiable, pas
seulement décrites en langage naturel. Quatre propriétés reviennent
systématiquement dans l'analyse de protocoles à vocation forensique
ou juridique~:

\begin{tableaucours}{p{3.5cm}p{3.0cm}p{4.5cm}}{%
    \tableauHeader{Propriété} &
    \tableauHeader{Lemme typique} &
    \tableauHeader{Pertinence forensique}
}
Non-répudiation & \texttt{lemma\_nonRep\_origin} &
    Un signataire ne peut nier avoir produit une signature,
    fondement de la valeur probatoire d'eIDAS
    (Ch.~11~(UE~SEC4031)) \\
Confidentialité & \texttt{lemma\_confidentiality} &
    Aucune information sur le message n'est exposée au-delà de ce
    qui est strictement nécessaire, correspond à $\mathsf{Conf}$
    (Ch.~\ref{chap:CRO}) \\
Authenticité & \texttt{lemma\_unforgeability} &
    Une signature ne peut être forgée sans la clé privée
    correspondante, fondement de la chaîne de custody
    numérique (Ch.~8~(UE~SEC4031)) \\
Intégrité (engagement) & \texttt{lemma\_binding} &
    Une preuve ne peut être réinterprétée \textit{a posteriori}
    pour correspondre à un message différent, analogue au
    hash SHA-256 garantissant qu'une image E01 n'a pas été
    altérée (Ch.~8~(UE~SEC4031)) \\
\end{tableaucours}
\vspace{-0.8em}
\begin{center}{\small\itshape Tableau~31.1. Quatre propriétés de
sécurité fondamentales et leur formalisation typique}\end{center}

% ============================================================================
\section{Outils d'Analyse Formelle}
\label{sec:outils-formels}
% ============================================================================

\subsection{Le \textit{prover} Tamarin}

\termecle{Tamarin Prover}\index{Tamarin Prover} est l'outil de
référence pour les preuves symboliques de protocoles. Il modélise
le protocole et les capacités de l'adversaire (typiquement,
Dolev-Yao) via des \textbf{règles de réécriture}, et permet de
vérifier automatiquement des propriétés de sécurité exprimées en
\textbf{logique temporelle}, c'est-à-dire des propriétés qui
portent non seulement sur les états du système mais sur la
\textbf{séquence} des événements (``si~$A$ se produit, alors~$B$
s'est nécessairement produit avant'').

\begin{boxoutils}{Principe de fonctionnement de Tamarin~(vue conceptuelle)}
\begin{lstlisting}[language=bash_forensique]
# STRUCTURE TYPE D'UN MODELE TAMARIN (illustration pedagogique,
# pas un fichier executable -- la syntaxe Tamarin reelle est plus
# riche ; se referer a la documentation officielle du projet
# pour toute utilisation effective)

# 1. Declaration des regles de reecriture (le protocole)
rule Register_Key:
    [ Fr(~ltk) ]
    --[ KeyGen($A, ~ltk) ]->
    [ !Ltk($A, ~ltk), !Pk($A, pk(~ltk)) ]

rule Sign_Message:
    [ !Ltk($A, ~ltk), In(m) ]
    --[ Signed($A, m) ]->
    [ Out(sign(m, ~ltk)) ]

# 2. Declaration des capacites de l'adversaire (Dolev-Yao
#    standard, generalement implicite dans Tamarin)

# 3. Formalisation des lemmes (proprietes a verifier)
lemma non_repudiation:
    "All A m #i. Signed(A, m) @ i ==>
        (Ex #j. KeyGen(A, ltk) @ j & j < i)"

# 4. Execution : Tamarin explore l'espace des etats accessibles
#    a l'adversaire et tente soit de PROUVER le lemme (preuve
#    valide pour tous les etats), soit de trouver une TRACE
#    D'ATTAQUE qui le viole
\end{lstlisting}
\end{boxoutils}

\begin{boiteRemarque}{L'état ``non prouvé'' n'est pas une faille}
Un résultat fréquent en pratique~: Tamarin (ou ProVerif) termine
sur un état \textbf{``non prouvé''} (\textit{inconclusive}), ni
preuve de validité, ni trace d'attaque trouvée. Ce résultat
\textbf{n'indique pas une faille du protocole}, c'est une
limitation bien documentée de ces outils, due à
l'\textbf{explosion de l'espace d'états}~: pour des protocoles
complexes combinant plusieurs primitives, l'espace des
comportements possibles de l'adversaire devient trop vaste pour
une exploration automatique exhaustive en temps fini. Un résultat
``non prouvé'' signale la nécessité de~: (a)~preuves manuelles
complémentaires ciblées sur les parties non couvertes, ou
(b)~une simplification du modèle (au prix d'une perte de
généralité), ou (c)~un découpage du protocole en sous-composants
analysés séparément. \textbf{Cette distinction, absence de
preuve $\neq$ preuve d'absence, est rigoureusement la même
que celle pratiquée en investigation forensique}~: l'absence
d'une preuve dans un dossier ne démontre jamais l'innocence, elle
signale seulement une question restée sans réponse
(Ch.~7~(UE~SEC4031), distinction constatation/analyse).
\end{boiteRemarque}

\subsection{Panorama comparatif des outils}

\begin{tableaucours}{p{2.5cm}p{3.5cm}p{4.5cm}}{%
    \tableauHeader{Outil} &
    \tableauHeader{Type de preuve} &
    \tableauHeader{Application principale}
}
\textbf{Tamarin} & Preuve symbolique &
    Protocoles complexes, propriétés temporelles (séquences
    d'événements, ordre causal entre actions) \\
\textbf{ProVerif} & Vérification automatique &
    Protocoles plus simples, propriétés d'équivalence
    (indistinguabilité de deux exécutions) \\
\textbf{CryptoVerif} & Preuve computationnelle &
    Preuves dans le modèle standard (au-delà du modèle
    symbolique~: tient compte des probabilités d'attaque,
    pas seulement de leur possibilité logique) \\
\end{tableaucours}
\vspace{-0.8em}
\begin{center}{\small\itshape Tableau~31.2. Comparatif des principaux
outils d'analyse formelle}\end{center}

\begin{boiteRemarque}{Symbolique vs computationnel~: une distinction
qui rejoint le Ch.~\ref{chap:CRO}}
Tamarin et ProVerif opèrent dans le \textbf{modèle symbolique}
(Dolev-Yao~: les primitives sont des boîtes noires parfaites).
CryptoVerif opère dans le \textbf{modèle computationnel}, plus
proche de l'approche information-théorique du cadre CLO
(Ch.~\ref{chap:CRO})~: les primitives y sont modélisées avec leurs
véritables propriétés probabilistes (négligeabilité, avantage de
l'adversaire), pas comme des oracles parfaits. Un protocole
``prouvé sûr'' dans le modèle symbolique offre une garantie plus
faible qu'un protocole ``prouvé sûr'' dans le modèle computationnel~:
le premier suppose que les primitives sous-jacentes sont parfaites
(hypothèse simplificatrice), le second quantifie explicitement le
risque résiduel, exactement la démarche des Théorèmes~1 à~3 du
Ch.~\ref{chap:CRO}, qui bornent $\mathsf{Conf}$, $\mathsf{Rel}$,
$\mathsf{Opp}$ par des quantités explicitement probabilistes
($\mathrm{negl}(\lambda)$), pas par de simples impossibilités
logiques.
\end{boiteRemarque}

% ============================================================================
\section{Méthodologie d'Audit en Cinq Étapes}
\label{sec:methodologie-audit}
% ============================================================================

Cette méthodologie systématique, standard dans la pratique
professionnelle de l'audit de protocoles cryptographiques, guide
l'analyse de tout protocole, qu'il s'agisse d'un nouveau schéma
proposé ou d'un mécanisme existant que l'investigateur doit évaluer
en contexte forensique (Ch.~\ref{chap:rapport}, témoignage
d'expert sur la fiabilité d'un mécanisme contesté).

\begin{tableaucours}{p{2.3cm}p{8.7cm}}{%
    \tableauHeader{Étape} &
    \tableauHeader{Contenu}
}
\textbf{1. Compréhension} & Analyse approfondie du \textit{whitepaper},
    de la documentation et, si disponible, du code source.
    Identification des acteurs, des messages échangés, des
    primitives cryptographiques employées et des objectifs de
    sécurité annoncés. \\
\textbf{2. Modélisation} & Définition explicite du modèle de menace
    (typiquement Dolev-Yao, Section~\ref{sec:modelisation-menaces})
    et formalisation des propriétés de sécurité attendues sous
    forme de lemmes vérifiables (Tableau~31.1). \\
\textbf{3. Analyse manuelle} & Recherche ciblée de vulnérabilités
    connues~: attaque par rejeu de \textit{nonce}, faiblesse du
    générateur de nombres aléatoires, mauvaise gestion de l'état
    des sessions, erreurs de composition entre primitives
    individuellement sûres (cf.~Tableau~30.1, faille ``contre le
    protocole''). \\
\textbf{4. Analyse automatisée} & Implémentation du modèle dans
    un outil tel que Tamarin ou ProVerif (Section~\ref{sec:outils-formels})
    pour une vérification exhaustive des propriétés contre un
    adversaire actif, dans les limites de l'explosion de l'espace
    d'états. \\
\textbf{5. Test d'implémentation} & Si une implémentation concrète
    est disponible~: tests pratiques (\textit{fuzzing} des entrées,
    analyse dynamique, tests de performance sous charge) et analyse
    statique du code source pour détecter des bugs d'implémentation
    (Tableau~30.1, faille ``d'implémentation'', distincte d'une
    faille conceptuelle). \\
\end{tableaucours}
\vspace{-0.8em}
\begin{center}{\small\itshape Tableau~31.3. Méthodologie d'audit
en cinq étapes}\end{center}

\begin{boxscenario}{Application de la méthodologie à un cas du cours~:
audit d'un mécanisme de signature de rapport}
Supposons qu'un nouveau mécanisme de signature pour les rapports
d'expertise forensique camerounais (Ch.~\ref{chap:rapport}) vous
soit soumis pour audit.

\textbf{Étape~1 (Compréhension)}~: lecture de la spécification~:
le mécanisme combine une signature ECDSA (Ch.~\ref{chap:prim-crypt})
avec un horodatage qualifié.

\textbf{Étape~2 (Modélisation)}~: modèle Dolev-Yao standard.
Lemmes à vérifier~: \texttt{non\_repudiation} (l'expert ne peut nier
avoir signé), \texttt{temporal\_integrity} (l'horodatage ne peut
être falsifié indépendamment de la signature).

\textbf{Étape~3 (Analyse manuelle)}~: vérifier si le mécanisme est
vulnérable à une attaque par rejeu, un rapport signé pourrait-il
être présenté une seconde fois, dans une affaire différente, comme
preuve falsifiée~? Cela dépend de la présence d'un identifiant de
session unique dans le message signé.

\textbf{Étape~4 (Analyse automatisée)}~: modélisation dans
Tamarin~; si le résultat est ``non prouvé'' pour la propriété de
non-rejeu, cela signale, conformément à la
Section~\ref{sec:outils-formels}, la nécessité d'une preuve
manuelle complémentaire, pas nécessairement une faille avérée.

\textbf{Étape~5 (Test d'implémentation)}~: si une implémentation
logicielle existe pour ce mécanisme, fuzzer les champs d'entrée
(date, identifiant d'affaire) pour détecter d'éventuels bugs de
validation.
\end{boxscenario}

% ============================================================================
\section{Rigueur sur les Travaux en Cours~: une Réserve Méthodologique}
\label{sec:reserve-zk-nr}
% ============================================================================

\begin{boiteRemarque}{Cohérence avec le Chapitre~\ref{chap:cryptanal}}
Le Ch.~\ref{chap:cryptanal} a mentionné le protocole de recherche
ZK-NR (\textit{Zero-Knowledge Non-Repudiation}), développé par
l'auteur de ce cours, en signalant explicitement qu'aucun article
complet correspondant n'était disponible pour vérification dans le
cadre de ce projet. La même réserve s'applique ici~: toute
affirmation précise sur le nombre de règles symboliques, de lemmes
formalisés, ou sur le statut exact (prouvé / non prouvé / réfuté)
d'une vérification Tamarin appliquée à ce protocole spécifique
\textbf{ne peut être incluse dans ce cours sans accès à une
publication ou un rapport technique vérifiable}.

\textbf{Ce qui peut être affirmé en toute rigueur}~: la méthodologie
en cinq étapes présentée dans ce chapitre (Section~\ref{sec:methodologie-audit})
constitue le cadre approprié pour conduire, le moment venu, l'audit
formel de tout protocole combinant plusieurs primitives, y compris
un protocole tel que ZK-NR, combinant engagements, preuves à
divulgation nulle et signatures à seuil. L'application effective
de cette méthodologie, et la publication de ses résultats
(prouvés, réfutés, ou non conclusifs), relève d'un travail de
recherche distinct de ce cours.

Cette transparence sur les limites de ce qui peut être affirmé,
ici comme au Ch.~\ref{chap:cryptanal}, illustre elle-même le
principe central de ce chapitre~: \textbf{distinguer rigoureusement
ce qui est prouvé, ce qui est en cours d'investigation, et ce qui
demeure une question ouverte}, exactement la discipline exigée
d'un rapport forensique (Ch.~\ref{chap:rapport}).
\end{boiteRemarque}

% ============================================================================
\section{De la Preuve de Protocole à la Preuve Forensique}
\label{sec:preuve-protocole-forensique}
% ============================================================================

\begin{boxCRO}
\textbf{Analyse CRO. Vérification formelle et opposabilité}

La vérification formelle d'un protocole, quand elle aboutit à une
preuve complète, constitue un argument d'opposabilité d'un type
particulier~: elle ne dépend pas de la confiance accordée à un
expert individuel, mais de la \textbf{vérifiabilité mécanique et
reproductible} de la preuve elle-même par n'importe quel tiers
disposant de l'outil utilisé.

$\mathcal{O}$~: une preuve Tamarin complète, soumise au tribunal,
peut en principe être \textbf{rejouée indépendamment} par un
expert de la partie adverse, un niveau de vérifiabilité supérieur
à un simple rapport d'expertise narratif (Ch.~\ref{chap:rapport}).

$\mathcal{R}$~: cette garantie est cependant \textbf{relative au
modèle de menace posé} (Dolev-Yao, Section~\ref{sec:modelisation-menaces})~:
une preuve formelle de sécurité ne dit rien sur la résistance aux
attaques par canal auxiliaire, aux failles d'implémentation, ou
à la cryptanalyse des primitives elles-mêmes (Ch.~\ref{chap:cryptanal}).

\cro{$\sim$}{$\uparrow$ dans le modèle posé~; limité hors modèle}{$\uparrow\uparrow$ rejouabilité indépendante}

\textbf{Implication pour le témoignage d'expert}
(Ch.~\ref{chap:rapport})~: un expert présentant une preuve formelle
doit, comme pour toute conclusion technique, \textbf{énoncer
explicitement le modèle de menace sous lequel la preuve est
valide}, exactement la discipline de la section
``Limitations'' déjà pratiquée dans le modèle de rapport du
Ch.~\ref{chap:rapport}.
\end{boxCRO}

\separateur

\begin{boxresume}
\textbf{Ce qu'il faut retenir de ce chapitre~:}
\begin{itemize}
    \item Le \textbf{modèle Dolev-Yao} sépare délibérément l'analyse
          logique du protocole (son objet) de la cryptanalyse des
          primitives (objet du Ch.~\ref{chap:cryptanal})~: l'adversaire
          contrôle le réseau mais ne peut casser les primitives
          par force brute.

    \item \textbf{Quatre propriétés fondamentales} à formaliser~:
          non-répudiation, confidentialité, authenticité, intégrité
          (engagement), chacune directement reliable à un
          mécanisme forensique déjà rencontré dans ce cours.

    \item \textbf{Tamarin Prover}~: outil de référence pour les
          preuves symboliques avec propriétés temporelles. Un
          résultat ``non prouvé'' \textbf{n'est pas une faille}~:
          c'est une limitation par explosion de l'espace d'états,
          signalant un besoin de preuve manuelle complémentaire.

    \item \textbf{Trois outils complémentaires}~: Tamarin
          (symbolique, temporel), ProVerif (automatique,
          équivalence), CryptoVerif (computationnel, quantifie le
          risque résiduel plutôt que de poser des primitives
          parfaites), ce dernier rejoignant l'esprit
          probabiliste du cadre CLO (Ch.~\ref{chap:CRO}).

    \item \textbf{Méthodologie en cinq étapes}~: compréhension,
          modélisation, analyse manuelle, analyse automatisée, test
          d'implémentation, applicable à tout protocole, y
          compris en contexte d'audit forensique d'un mécanisme
          contesté.

    \item \textbf{Rigueur sur les affirmations non vérifiables}~:
          ce chapitre s'est volontairement abstenu de détails
          précis sur le protocole de recherche ZK-NR, faute
          d'article vérifiable, cohérence méthodologique directe
          avec la réserve posée au Ch.~\ref{chap:cryptanal}.

    \item \textbf{Vérification formelle et opposabilité}~: une
          preuve formelle complète et rejouable constitue un
          argument d'opposabilité fort, mais \textbf{toujours
          relatif au modèle de menace posé}, à énoncer
          explicitement dans tout témoignage d'expert
          (Ch.~\ref{chap:rapport}).
\end{itemize}
\end{boxresume}

\bigskip

\begin{boxexo}[title={\ding{45}\quad Exercice~31.1. Maîtrise du modèle Dolev-Yao \niveaubase}]
\begin{enumerate}[(a)]
    \item Le modèle Dolev-Yao suppose que l'adversaire ne peut pas
          casser les primitives cryptographiques par force brute.
          En quoi cette hypothèse diffère-t-elle de l'approche du
          cadre CLO (Ch.~\ref{chap:CRO}), qui borne explicitement
          un terme $\mathrm{negl}(\lambda)$ plutôt que de le
          supposer nul~?
    \item Pour chacune des quatre propriétés du Tableau~31.1,
          identifiez le mécanisme forensique de ce cours qui
          l'illustre le mieux, en justifiant.
    \item Un protocole de vote électronique permet à un adversaire
          de rejouer un bulletin déjà soumis. À quelle étape de la
          méthodologie d'audit (Tableau~31.3) cette vulnérabilité
          serait-elle le plus probablement détectée en premier~?
          Justifiez.
\end{enumerate}
\end{boxexo}

\begin{boxexo}[title={\ding{45}\quad Exercice~31.2. Interprétation d'un résultat de vérification \niveaumoyen}]
Un outil Tamarin appliqué à un protocole de signature de rapport
forensique (sur le modèle du \texttt{boxscenario} de ce chapitre)
renvoie un résultat ``non prouvé'' pour le lemme
\texttt{non\_repudiation}.
\begin{enumerate}[(a)]
    \item Pourquoi serait-il méthodologiquement incorrect de
          conclure directement que le protocole est vulnérable~?
    \item Quelles sont les trois pistes mentionnées dans ce
          chapitre (Section~\ref{sec:outils-formels}) pour
          progresser face à un tel résultat~? Laquelle
          privilégieriez-vous pour un protocole destiné à un usage
          judiciaire à haute criticité, et pourquoi~?
    \item En vous appuyant sur l'analogie posée dans ce chapitre
          (absence de preuve $\neq$ preuve d'absence), rédigez un
          paragraphe expliquant comment cette même règle
          s'applique à la rédaction d'un rapport forensique
          (Ch.~\ref{chap:rapport})~: que doit faire un investigateur
          qui n'a pas trouvé de preuve d'une hypothèse, sans pouvoir
          non plus la réfuter formellement~?
\end{enumerate}
\end{boxexo}

\begin{boxexo}[title={\ding{45}\quad Exercice~31.3. Audit complet d'un protocole proposé \niveauexpert}]
Un développeur camerounais propose un nouveau protocole de
``scellement numérique'' pour les pièces à conviction électroniques,
destiné à être utilisé par les OPJ (Ch.~8~(UE~SEC4031)). Le
protocole combine~: un hachage SHA-256 de l'evidence, une signature
ECDSA de l'OPJ, et un horodatage via une autorité tierce.
\begin{enumerate}[(a)]
    \item Appliquez l'Étape~1 (Compréhension) de la méthodologie~:
          identifiez les acteurs, les messages échangés, les
          primitives utilisées et les trois objectifs de sécurité
          que ce protocole devrait viser (référez-vous au
          Ch.~8~(UE~SEC4031)).
    \item Appliquez l'Étape~2 (Modélisation)~: formalisez au moins
          deux lemmes de sécurité pertinents pour ce protocole,
          sur le modèle du Tableau~31.1.
    \item Appliquez l'Étape~3 (Analyse manuelle)~: identifiez au
          moins une vulnérabilité potentielle de composition entre
          les trois primitives (hachage~+ signature~+ horodatage), par exemple, que se passe-t-il si l'horodatage est
          obtenu \emph{avant} la signature plutôt qu'après~?
    \item Pour les Étapes~4 et~5, décrivez (sans les exécuter
          réellement) ce que vous chercheriez à vérifier
          automatiquement, et quels tests d'implémentation
          seraient prioritaires si ce protocole devait être
          déployé dans les unités PNC.
    \item Rédigez un court paragraphe de synthèse destiné à
          l'ANTIC, présentant votre niveau de confiance dans ce
          protocole à l'issue de cet audit partiel, en respectant
          scrupuleusement la distinction entre ce que votre analyse
          établit et ce qui resterait à vérifier avant tout
          déploiement opérationnel.
\end{enumerate}
\end{boxexo}

\bigskip

\begin{boxinfo}
\textbf{Transition.} La Partie~VI a parcouru l'ensemble du socle
théorique de ce cours~: formalisation du Trilemme CRO
(Ch.~\ref{chap:CRO}), application aux primitives réelles
(Ch.~\ref{chap:prim-crypt}), conception et cryptanalyse
(Ch.~\ref{chap:cryptanal}), et enfin l'outillage de la preuve
formelle (ce chapitre). Le chapitre suivant et dernier de ce cours
propose un \textbf{cas pratique intégrateur} mobilisant l'ensemble
de ces outils théoriques sur une situation forensique complète,
clôturant ainsi le parcours entamé au Prologue.
\end{boxinfo}